% GENERATED FILE. Edit canonical lessons, then run npm run generate:books.
% canonical-source-hash: fnv1a64:e11ec0221da688be
% canonical-lessons: BN-C135-bhije, BN-C135-shukno, BN-C135-bhora, BN-C135-druto, BN-C135-dhir

\chapter{Wet, Dry, Full, Fast, Slow}
\label{ch:bn-wet-dry-full-fast-slow}
\hlchaptermodality{135}

\begin{chapteropening}
\textbf{By the end of this chapter:} \emph{I can say wet, dry, full, fast and slow.}

Complete the last lesson of chapter 135: i can say wet, dry, full, fast and slow.
\end{chapteropening}

\section[\texorpdfstring{\bn{ভিজে} (bhije)}{ভিজে (bhije)}]{\bn{ভিজে} (bhije) — wet}

\label{lesson:BN-C135-bhije}



\begin{quote}
Before the new one: say the Bengali for hard, then the Bengali for soft.
\end{quote}

\subsection*{You'll want to know: \bn{ভিজে}}
\textbf{\bn{ভিজে}} — \emph{bhije} — \textquotedblleft{}wet\textquotedblright{}.

\begin{grammarlens}[title={What you've built}]
One more word for this chapter.
\end{grammarlens}

\subsection*{Your turn}
\begin{itemize}
\raggedright
  \item \emph{Say it:} \emph{bhije}
  \item \emph{Say it:} \emph{bhije}, once more
  \item \emph{Say it:} say \emph{bhije}
  \item \emph{Recall:} say the Bengali for hard, then the Bengali for soft, then say \emph{bhije} again
\end{itemize}

\begin{tcolorbox}[breakable,colback=teal!4,colframe=teal!35!black,title={Before you move on}]
What does \bn{ভিজে} mean? (\textquotedblleft{}Wet\textquotedblright{}.) Say it once more.
\end{tcolorbox}
\section[\texorpdfstring{\bn{শুকনো} (shukno)}{শুকনো (shukno)}]{\bn{শুকনো} (shukno) — dry}

\label{lesson:BN-C135-shukno}



\begin{quote}
Before the new one: say the Bengali for soft, then the Bengali for wet.
\end{quote}

\subsection*{You'll want to know: \bn{শুকনো}}
\textbf{\bn{শুকনো}} — \emph{shukno} — \textquotedblleft{}dry\textquotedblright{}.

\begin{grammarlens}[title={What you've built}]
One more word for this chapter.
\end{grammarlens}

\subsection*{Your turn}
\begin{itemize}
\raggedright
  \item \emph{Say it:} \emph{shukno}
  \item \emph{Say it:} \emph{shukno}, once more
  \item \emph{Say it:} say \emph{shukno}
  \item \emph{Recall:} say the Bengali for soft, then the Bengali for wet, then say \emph{shukno} again
\end{itemize}

\begin{tcolorbox}[breakable,colback=teal!4,colframe=teal!35!black,title={Before you move on}]
What does \bn{শুকনো} mean? (\textquotedblleft{}Dry\textquotedblright{}.) Say it once more.
\end{tcolorbox}
\section[\texorpdfstring{\bn{ভরা} (bhôrā)}{ভরা (bhôrā)}]{\bn{ভরা} (bhôrā) — full}

\label{lesson:BN-C135-bhora}



\begin{quote}
Before the new one: say the Bengali for wet, then the Bengali for dry.
\end{quote}

\subsection*{You'll want to know: \bn{ভরা}}
\textbf{\bn{ভরা}} — \emph{bhôrā} — \textquotedblleft{}full\textquotedblright{}.

\begin{grammarlens}[title={What you've built}]
One more word for this chapter.
\end{grammarlens}

\subsection*{Your turn}
\begin{itemize}
\raggedright
  \item \emph{Say it:} \emph{bhôrā}
  \item \emph{Say it:} \emph{bhôrā}, once more
  \item \emph{Say it:} say \emph{bhôrā}
  \item \emph{Recall:} say the Bengali for wet, then the Bengali for dry, then say \emph{bhôrā} again
\end{itemize}

\begin{tcolorbox}[breakable,colback=teal!4,colframe=teal!35!black,title={Before you move on}]
What does \bn{ভরা} mean? (\textquotedblleft{}Full\textquotedblright{}.) Say it once more.
\end{tcolorbox}
\section[\texorpdfstring{\bn{দ্রুত} (drutô)}{দ্রুত (drutô)}]{\bn{দ্রুত} (drutô) — fast}

\label{lesson:BN-C135-druto}



\begin{quote}
Before the new one: say the Bengali for dry, then the Bengali for full.
\end{quote}

\subsection*{You'll want to know: \bn{দ্রুত}}
\textbf{\bn{দ্রুত}} — \emph{drutô} — \textquotedblleft{}fast\textquotedblright{}.

\begin{grammarlens}[title={What you've built}]
One more word for this chapter.
\end{grammarlens}

\subsection*{Your turn}
\begin{itemize}
\raggedright
  \item \emph{Say it:} \emph{drutô}
  \item \emph{Say it:} \emph{drutô}, once more
  \item \emph{Say it:} say \emph{drutô}
  \item \emph{Recall:} say the Bengali for dry, then the Bengali for full, then say \emph{drutô} again
\end{itemize}

\begin{tcolorbox}[breakable,colback=teal!4,colframe=teal!35!black,title={Before you move on}]
What does \bn{দ্রুত} mean? (\textquotedblleft{}Fast\textquotedblright{}.) Say it once more.
\end{tcolorbox}
\section[\texorpdfstring{\bn{ধীর} (dhīr)}{ধীর (dhīr)}]{\bn{ধীর} (dhīr) — slow}

\label{lesson:BN-C135-dhir}



\begin{quote}
Before the new one: say the Bengali for full, then the Bengali for fast.
\end{quote}

\subsection*{You'll want to know: \bn{ধীর}}
\textbf{\bn{ধীর}} — \emph{dhīr} — \textquotedblleft{}slow\textquotedblright{}.

\begin{grammarlens}[title={What you've built}]
One more word for this chapter.
\end{grammarlens}

\subsection*{Your turn}
\begin{itemize}
\raggedright
  \item \emph{Say it:} \emph{dhīr}
  \item \emph{Say it:} \emph{dhīr}, once more
  \item \emph{Say it:} say \emph{dhīr}
  \item \emph{Recall:} say the Bengali for full, then the Bengali for fast, then say \emph{dhīr} again
\end{itemize}

\begin{tcolorbox}[breakable,colback=teal!4,colframe=teal!35!black,title={Before you move on}]
What does \bn{ধীর} mean? (\textquotedblleft{}Slow\textquotedblright{}.) Say it once more.
\end{tcolorbox}
